\chapter{Multilinear Algebra}
\section{Tensor Product}

\quad We fix a commutative unitary ring $K$.
For a natural number $n\in \NN_{\ge 1}$, $V_1,\cdots, V_n$, $W$ collection of $K$-modules,
we denote by $\mathrm{Hom}_{K}^{(n)} \left(V_1,\cdots,V_n;W\right)$ the set of $n$-linear mappings from $V_1\times\cdots\times V_n$ to $W$.

\begin{theoremenv}\label{thm:4.1.1}
    Let $n$ be a natural number, $V_1,\cdots, V_n$ be $K$-modules.
    The functor $F$ from $\mathbf{Mod}_{K}$ to $\mathbf{Set}$ that sends any $K$-module $W$ to the set of $n$-linear mappings from $V_1\times\cdots\times V_n$ to $W$ is representable.
    $$\begin{array}{rrcl}
        F: & \mathbf{Mod}_{K} & \longrightarrow & \mathbf{Set}\\
        & W & \longmapsto & \mathrm{Hom}_{K}^{(n)} \left(V_1,\cdots,V_n;W\right).
    \end{array}$$
\end{theoremenv}
\begin{proofenv}
    We have seen that, for any $K$-module $W$, the mapping
    $$ W^{V_1\times \cdots \times V_n} \rightarrow \mathrm{Hom}_{K} \left(K^{\oplus (V_1\times \cdots \times V_n)}, W\right)$$
    sending any $s \in W^{V_1\times \cdots \times V_n}$ to 
    $$\begin{array}{rrcl}
        f_{s} : & K^{\oplus (V_1\times \cdots \times V_n)} & \longrightarrow & W\\
        & \left(a_{\bm{x}}\right)_{\bm{x}\in V_1\times \cdots \times V_n} & \longmapsto &\displaystyle \sum_{\bm{x}\in V_1\times \cdots \times V_n} a_{\bm{x}} s(\bm{x})
    \end{array}$$
    is a bijection (example about $G: \mathbf{Mod}_{K} \rightarrow \mathbf{Set}$, $W \mapsto W^I$ is representable by $K^{\oplus I}$.)
    Let $\left(e_{\bm{x}}\right)_{\bm{x}\in V_1 \times \cdots \times V_n}$ be the canonical basis of $K^{\oplus (V_1\times \cdots \times V_n)}$ which means that $e_{\bm{x}}$ sends $\bm{x}$ to $1$ and any $\bm{y}\in V_{1}\times \cdots \times V_n\setminus\{\bm{x}\}$ to $0$.
    Let $M$ be the $K$-submodule of $K^{\oplus V_1\times \cdots \times V_n}$ generated by elements of the form
    \begin{equation*}
        \begin{aligned}
            & e\left(x_1,\cdots,x_{i-1}, ax_i' + x_i'',x_{i+1},\cdots,x_n\right)\\
            & - a e\left(x_1,\cdots,x_{i-1}, x_i',x_{i+1},\cdots,x_n\right)\\
            & - b e\left(x_1,\cdots,x_{i-1}, x_i'',x_{i+1},\cdots,x_n\right)
        \end{aligned} \tag{$*$}
    \end{equation*}
    where $a,b\in K$, $x_i', x_i'' \in V_i$.
    $f_s$ associated to $s \in W^{V_1\times V_n}$ sends ($*$) to  
        \begin{equation*}
        \begin{aligned}
            & s\left(x_1,\cdots,x_{i-1}, ax_i' + x_i'',x_{i+1},\cdots,x_n\right)\\
            & - a s\left(x_1,\cdots,x_{i-1}, x_i',x_{i+1},\cdots,x_n\right)\\
            & - b s\left(x_1,\cdots,x_{i-1}, x_i'',x_{i+1},\cdots,x_n\right)
        \end{aligned}
    \end{equation*}
    and therefore $s$ is linear if and only if $f_s$ vanishes on $M$.
    In this case, $f_s$ is induced homomorphism from
    $$K^{\oplus (V_1\times \cdots \times V_n)}/M \eqcolon V_1\otimes \cdots \otimes V_n.$$
    We conclude that $F$ is representable by the object $V_1\otimes \cdots \otimes V_n$.
\end{proofenv}

\begin{definitionenv}
    Let $n\in \NN$, $V_1,\cdots, V_n$ be $K$-modules.
    The $K$-module $V_1\otimes \cdots \otimes V_n$ is called the \textbf{tensor product} of $V_1,\cdots, V_n$.
    For any $(x_1,\cdots,x_n) \in V_1\times \cdots \times V_n$, we denote $x_1\otimes \cdots \otimes x_n$ the image of $(x_1,\cdots,x_n)$ by the universal $n$-linear mapping
    $$ V_1\times \cdots \times V_n \rightarrow V_1\otimes \cdots \otimes V_n.$$
    Note that for any $K$-module $W$ and any $n$-linear mapping $f: V_1\times \cdots \times V_n \rightarrow W$, there exists a \textit{unique} homomorphism of $K$-modules such that the following diagram commutes
    \begin{center}
    \begin{tikzcd}
    V_1 \times \cdots \times V_n \arrow[d, ""] \arrow[dr, "f"]& \\
    V_1 \otimes \cdots \otimes V_n \arrow[r, swap, "\exists! \tilde{f}"]& W 
    \end{tikzcd}
    \end{center}

Sometimes $V_1 \otimes \cdots \otimes V_n$ is denoted by $V_1 \otimes_{K} \cdots \otimes_{K} V_n$.
\end{definitionenv}

\begin{propositionenv}
    Let $n\in \NN_{\ge 2}$, $V_1, \cdots, V_n$ be $K$-modules. 
    Let $i \in \{1,\cdots, n-1\}$ and $j \in \{0, \cdots, n-i\}$.
    There is an isomorphism of $K$-modules
    \begin{equation*}
        \begin{array}{rcl}
            V_1 \otimes \cdots \otimes V_n & \rightarrow & V_1 \otimes \cdots \otimes V_{i-1} \otimes \left(V_i \otimes \cdots \otimes V_{i+j}\right) \otimes V_{i+j+1} \otimes \cdots \otimes V_n\\
            x_1 \otimes \cdots \otimes x_n & \mapsto & x_1 \otimes \cdots \otimes x_{i-1} \otimes (x_i \otimes \cdots \otimes x_{i+j}) \otimes x_{i+j+1} \otimes \cdots \otimes x_n.
        \end{array}
    \end{equation*}
\end{propositionenv}
\begin{proofenv}
    The mapping
    \begin{equation*}
        \begin{array}{rcl}
            V_1 \times \cdots \times V_n & \rightarrow & V_1 \otimes \times \otimes V_{i-1} \otimes \left(V_i \otimes \cdots \otimes V_{i+j}\right) \otimes V_{i+j+1} \otimes \cdots \otimes V_n\\
            (x_1 \times \cdots \times x_n) & \mapsto & x_1 \otimes \cdots \otimes x_{i-1} \otimes (x_i \otimes \cdots \otimes x_{i+j}) \otimes x_{i+j+1} \otimes \cdots \otimes x_n
        \end{array}
    \end{equation*}    
    is linear, hence induces a $K$-module homomorphism
    \begin{equation*}
        \begin{array}{rcl}
            V_1 \otimes \cdots \otimes V_n & \rightarrow & V_1 \otimes \times \otimes V_{i-1} \otimes \left(V_i \otimes \cdots \otimes V_{i+j}\right) \otimes V_{i+j+1} \otimes \cdots \otimes V_n\\
            x_1 \otimes \cdots \otimes x_n & \mapsto & x_1 \otimes \cdots \otimes x_{i-1} \otimes (x_i \otimes \cdots \otimes x_{i+j}) \otimes x_{i+j+1} \otimes \cdots \otimes x_n.
        \end{array}
    \end{equation*}
    It remains to check that this is actually an isomorphism.
    To achieve it, it is enough to show that both modules represent the same functor.
    $$\begin{array}{rrcl}
        F: & \mathbf{Mod}_{K} & \rightarrow & \mathbf{Set}\\
        & W & \mapsto & \mathrm{Hom}_{K}^{(n)} \left(V_1\times\cdots\times V_n,W\right).
    \end{array}$$
    By theorem \ref{thm:4.1.1}, we have the following bijections functorial with respect to $W$:
    \[
    \begin{aligned}
    &\mathrm{Hom}_K(V_1 \otimes \cdots \otimes V_n; W) \\
    &\rightarrow \mathrm{Hom}_K^{(n)}(V_1, \cdots, V_n; W) \\
    &\rightarrow \mathrm{Hom}_K^{(n-j)}(V_1, \cdots, V_{i-1}, V_i \otimes \cdots \otimes V_{i+j}, V_{i+j+1}, \cdots, V_n; W) \\
    &\rightarrow \mathrm{Hom}_K(V_1 \otimes \cdots \otimes V_{i-1} \otimes (V_i \otimes \cdots \otimes V_{i+j}) \otimes V_{i+j+1} \otimes \cdots \otimes V_n; W).
    \end{aligned}
    \]
\end{proofenv}

\section{Tensor Algebra}

\begin{definitionenv}
    Let $K$ be a commutative unitary ring and $E$ be a $K$-module.
    For any $n \in \NN_{\ge 1}$, we denote by $E^{\otimes n}$ the $K$-module
     $$\underset{n \text{ times}}{\underbrace{E\otimes \cdots\otimes E}}.$$
    By convention $E^{\otimes 0} = K.$
    Define
    $$ T(E) \coloneq \bigotimes_{n\in \NN} E^{\otimes n}.$$

    For any $n\in \NN$, we have a homomorphism of $K$-modules
    $$ E^{\otimes n} \rightarrow T(E)$$
    that sends any $\alpha_n \in E^{\otimes n}$ to the element of $T(E)$ whose $n$-th coordinate equals $\alpha_n$ and all the other coordinates equals zero.
    This mapping allows us to see $E^{\otimes n}$ as a $K$-submodule of $T(E)$.
    By abuse of notation, an element $\left(\alpha_n\right)_{n\in \NN}$ is written as
    $$ \sum_{n\in \NN} \alpha_n.$$
    We consider he following bilinear map:
    $$ \otimes : T(E) \times T(E) \rightarrow T(E)$$
    that sends $(x_1\otimes \cdots \otimes x_n, x_{n+1}\otimes \cdots \otimes x_{n+m}) \in E^{\otimes n} \times E^{\otimes m}$ to
    $$ x_1 \otimes \cdots \otimes x_n \otimes x_{n+1} \otimes \cdots \otimes x_{n+m} \in E^{\otimes (n+m)}.$$
    This equips $T(E)$ with a structure of $K$-algebra.
    It is called the \textbf{tensor algebra} of $E$.
    Sometimes we write $T_{K}(E)$ instead of $T(E)$.
\end{definitionenv}

\begin{theoremenv}
    Let $K$ be a commutative unitary ring, $E$ be a $K$-module.
    The functor from $\mathbf{Alg}_{K}$ to $\mathbf{Set}$ that sends any $K$-algebra $A$ to the set of $K$-module homomorphisms from $E$ to $A$ is representable by $T(E)$.
\end{theoremenv}
\begin{proofenv}
    It is suffices to check that for any $K$-algebra $A$, and any $K$-module homomorphism $f: E \rightarrow A$, $f$ extends in a unique way to a homomorphism of $K$-algebras $\lambda_f$ from $T(E)$ to $A$.

    \quad Assume that such $\lambda_f$ exists.
    Then for any $x_1 \otimes \cdots \otimes x_n \in E^{\otimes n}$,
    $$ \lambda_f \left(x_1 \otimes \cdots \otimes x_n\right) = f \left(x_1\right) \cdot \ldots \cdot f \left(x_n\right)$$
    which proves uniqueness of $\lambda_f$ since $T(E)$ is generated by the elements of the form $x_1 \otimes \cdots \otimes x_n$ for some $n\in \NN$ and $x_1, \cdots, x_n \in E$.

    \quad It remains to show existence of $\lambda_f$.
    First, consider the mapping:
    $$\begin{array}{rcl}
        E^n & \longrightarrow & A\\
        (x_1, \cdots, x_n) & \longmapsto & f(x_1) \cdot \ldots \cdot f(x_n).
    \end{array}$$
    The mapping is $n$-linear, hence induces a $K$-linear mapping $E^{\otimes n} \rightarrow A$ by universal property of $\otimes$.
    We can construct these $E^{\otimes n} \rightarrow A$ for any $n\in \NN$.
    The collection obtained this way gives us a $K$-linear mapping
    $$ \lambda_f : T(E) \longrightarrow A$$
    that sends $x_1 \otimes \cdots \otimes x_n$ to $f(x_1) \cdot \ldots \cdot f(x_n).$
    One has 
    $$\begin{aligned}
        &\lambda_f \left(x_1 \otimes \cdots \otimes x_n \otimes x_{n+1} \otimes \cdots \otimes x_{n+m}\right)\\
        = & f(x_1) \cdot \ldots \cdot f(x_n) \cdot f(x_{n+1}) \cdot \ldots \cdot f(x_{n+m})\\
        = & \left(f(x_1) \cdot \ldots \cdot f(x_n)\right) \cdot \left(f(x_{n+1}) \cdot \ldots \cdot f(x_{n+m})\right)\\
        =& \lambda_f \left(x_1 \otimes \cdots \otimes x_n\right) \cdot \lambda_f \left(x_{n+1} \otimes \cdots \otimes x_{n+m}\right).
    \end{aligned}$$
    So $\lambda_f$ is indeed a homomorphism of $K$-algebras.
\end{proofenv}

\begin{corollaryenv}
    Let $V$, $W$ be $K$-modules and $f : V\rightarrow W$ be a homomorphism of $K$-modules.
    Then $f$ induces a homomorphism of $K$-algebras 
    $$ \begin{array}{rrcl}
        T(f) : & T(V) & \longrightarrow & T(W)\\
        & x_1 \otimes \cdots \otimes x_n & \longmapsto & f(x_1) \otimes \cdots \otimes f(x_n).
    \end{array}$$
\end{corollaryenv}
\begin{proofenv}
    $V \overset{f}{\longrightarrow} W \overset{i}{\longrightarrow} T(W) $ is a homomorphism, where $i: W \longrightarrow T(W)$ is the inclusion.
    By the theorem, there exists a homomorphism of $K$-algebras $T(f) : T(V) \rightarrow T(W)$ such that the following diagram commutes:
    \begin{center}    \begin{tikzcd}
    V \arrow[r, "f"] \arrow[d, "i_V"] & W \arrow[d, "i_W"]\\
    T(V) \arrow[r, swap, "T(f)"] & T(W)
    \end{tikzcd}
    \end{center}
\end{proofenv}

\begin{remark}
    We call \textbf{graded $\bm{K}$-algebra} $A$ any $K$-algebra that can be written in the form
    $$ A = \bigoplus_{n\in \NN} A_n$$
    as a $K$-module, such that for any $n,m \in \NN$,
    $$ \forall a_n \in A_n, a_m \in A_m, a_n a_m \in A_{n+m}.$$
    
    \quad For example, if $V$ is a $K$-module, then the tensor algebra $T(V)$ is a graded $K$-algebra.
    Given a graded $K$-algebra
    $$ A = \bigoplus_{n\in \NN} A_n,$$
    for any $n \in \NN$, an element of $A_n$ is said to be \textbf{homogeneous of degree $\bm{n}$}.
    (In particular, the zero element of $A$ is homogeneous of degree $n$ for any $n\in \NN$.)
    Note that for any $a \in A$, $a$ can be written in a \text{unique} way as
    $$ a = \sum_{n\in \NN} a_n,\; a_n \in A_n.$$
    In particular, like in the case of tensor algebra, we consider $A_n$ as a $K$-submodule of $A$.
    The element $a_n$ is called the \textbf{homogeneous component of degree $\bm{n}$}.
    
    \quad Let $I$ be an ideal of $A$.
    We say that $I$ is a \textbf{homogeneous ideal of $\bm{A}$} if for any $a \in I$, the homogeneous components of $a$ belong to $I$ (this is equivalent to saying that $I$ is generated by homogeneous elements).
    Any homogeneous ideal $I$ in $A = \bigoplus_{n \in \NN} A_n$ can be written as
    $$ I = \bigoplus_{n\in \NN} I_n$$
    where $I_n$ is the $K$-submodule of $I$ consisting of all its homogeneous elements of degree $n$.
    The quotient ring $A/I$ can be written as
    $$ A/I = \bigoplus_{n\in \NN} A_n/I_n$$
    and therefore admits a graded $K$-algebra structure.
\end{remark}


\section{Symmetric Algebra}

\begin{definitionenv}
    For any set $X$, we denote by $\mathfrak{S}_X$ of all bijections from $X$ to $X$.
    The set $\mathfrak{S}_X$ equipped with the composition of mappings forms a group called the \textbf{symmetric group} of $X$.
    The elements of $\mathfrak{S}_X$ are called \textbf{permutations} of $X$.

    \quad If $x_1,\cdots,x_n$ are distinct elements of $X$, we denote by
    $$ (x_1\; x_2\; \cdots\; x_n)$$
    the mapping that sends $x_1$ to $x_2$, $x_2$ to $x_3$, $\cdots$, $x_{n-1}$ to $x_n$ and $x_n$ to $x_1$, and any other element of $X$ to itself.
    An element of this form is called an \textbf{$\bm{n}$-cycle}.
    A $2$-cycle is called a \textbf{transposition}.
    During the first semester, we have proved that for a finite set $X$, any permutation $\sigma \in \mathfrak{S}_n$ can be written as
    $$ \sigma = \tau_1 \cdots \tau_d$$
    where $\tau_i$, $1\le i\le d$ is an $n_i$-cycle on a subset $X_i \subseteq X$ for $X_1,\cdots,X_d$ being pairwise disjoint sets of cardinality $n_1,\cdots,n_d$ respectively.
\end{definitionenv}

\begin{remark}
    Let $n\in \NN_{\ge 2}$ be an integer.
    The symmetric group $\mathfrak{S}_{\{1,\cdots,n\}}$ is denoted by $\mathfrak{S}_{n}$.
    A transposition of the form $(i,i+1)$ is called an \textbf{adjacent transposition}.
    We have seen that any permutation can be written as a product of adjacent transpositions.
\end{remark}

\begin{theoremenv}
    Let $V$ be a $K$-module.
    The functor from $\mathbf{CAlg}_{K}$ to $\mathbf{Set}$ sending any commutative $K$-algebra $A$ to the set of $K$-module homomorphisms from $V$ to $A$ is representable.
\end{theoremenv}
\begin{proofenv}
    Let $I$ be the ideal of $T(V)$ generated by the elements of the form $x\otimes y - y\otimes x$, where $x,y \in V$.
    Since any permutation in $\mathfrak{S}_{n}$ can be written as a composition of adjacent transpositions, for any $x_1, \cdots, x_n \in V$ and any $\sigma \in \mathfrak{S}_n$, one has
    $$ x_1 \otimes \cdots \otimes x_n - x_{\sigma(1)} \otimes \cdots \otimes x_{\sigma(n)} \in I.$$
    Hence the equivalent classes of $x_1\otimes \cdots \otimes x_n$ and $x_{\sigma(1)} \otimes \cdots \otimes x_{\sigma(n)}$ in $S(V) \coloneq T(V)/I$ are the same.
    In particular, one can see that the quotient ring $S(V)$ is commutative.
    
    \quad Let $f: V \rightarrow A$ be a homomorphism of $K$-modules, where $A$ is a commutative $K$-algebra.
    It induces ni a unique way a homomorphism of $K$-algebras $\lambda_f : T(V) \rightarrow A$.
    In addition, for any $x,y \in V$,
    $$ \lambda_f(x\otimes y - y\otimes x) = f(x) \cdot f(y) - f(y) \cdot f(x) = 0.$$
    So $\lambda_f$ vanishes on $I$ and therefore induces a unique homomorphism of $K$-algebras from $S(V)$ to $A$ that extends $f$.
    We conclude that $S(V)$ represents the functor under consideration.
\end{proofenv}

\begin{definitionenv}
    Let $V$ be a $K$-module.
    The commutative $K$-algebra $\bm{S(V)}$ is called the \textbf{symmetric algebra} of $V$.
\end{definitionenv}

\begin{remark}
    The ideal $I$ such that $S(V) = T(V)/I$ is generated by homogeneous elements.
    Therefore $S(V)$ is a graded $K$-algebra.
    We denote by $S^n(V)$ the homogeneous component of $S(V)$ of degree $n$.
    $$ S(V) = \bigoplus_{n\in \NN} S^n(V).$$
    The composition law written multiplicatively in $S(V)$, for any $n\in \NN$, one has a natural surjective mapping
    $$ \begin{array}{rcl}
        V^{\otimes n} & \rightarrow & S^n(V)\\
        x_1 \otimes \cdots \otimes x_n & \mapsto & x_1 \cdots x_n.
    \end{array}$$
    Since $I_n$ is generated by elements of the form $x_1 \otimes \cdots \otimes x_n - x_{\sigma(1)} \otimes \cdots \otimes x_{\sigma(n)}$ 
    we obtain $S(V)$ is the quotient of $V^{\otimes n}$ by the linear action of the symmetric group $\mathfrak{S}_n$ as follows
    $$ \forall \sigma \in \mathfrak{S}_n, \sigma (x_1 \otimes \cdots \otimes x_n) = x_{\sigma(1)} \otimes \cdots \otimes x_{\sigma(n)}.$$
\end{remark}

\begin{exampleenv}
    Let $K$ be a commutative unitary ring and $V$ be a $K$-module of finite type, and $\{T_i\}_{i=1}^{d}$ be a basis of $V$ over $K$.
    For any commutative $K$-algebra $A$, the set of homomorphisms of $K$-modules from $V$ to $K$ is a functorial bijection with $A^{\{T_1,\cdots,T_d\}}$ (fancy way to say that a homomorphism of $K$-modules $V \rightarrow K$ is uniquely determined by $f(T_1), \cdots, f(T_d) \in A$).
    
    \quad We denote by $K[T_1,\cdots, T_d]$ the graded $K$-algebra $S(V)$, called the \textbf{polynomial algebra} of variables $T_1, \cdots, T_d$. 
\end{exampleenv}
\newpage

\section{Exterior Algebra}
In this section, $K$ stands for a commutative unitary ring.

\begin{definitionenv}
    Let $V$ be a $K$-module.
    We denote by $\bigwedge (V)$ the quotient $K$-algebra of $T(V)$ modulo the ideal generated by elements of the form $x\otimes x$ for $x \in V$.
    The homogeneous component of $\bigwedge (V)$ of degree $n$ is denoted by $\bigwedge^n (V)$.
    The composition law of $\bigwedge (V)$ is denoted by $\wedge$.
    If $n$ is a positive integer, and $x_1, \cdots, x_n$ are elements of $V$, then we have
    $$\begin{array}{rcl}
        V^{\otimes n} & \longrightarrow & \bigwedge^n (V),\\
        x_1 \otimes \cdots \otimes x_n & \longmapsto & x_1 \wedge \cdots \wedge x_n.
    \end{array}$$
    The $K$-algebra $\bigwedge (V)$ is called the \textbf{exterior algebra} of $V$.
\end{definitionenv}

\begin{propositionenv}
    Let $n$ be a natural number.
    There exists a homomorphism $\mathrm{sgn}$ from $\mathfrak{S}_n$ to $\{\pm 1\}$ that sends any transposition to $-1$.
\end{propositionenv}
\begin{proofenv}
    Consider the mapping
    $$\begin{array}{rrcl}
        \mathrm{sgn} : & \mathfrak{S}_n & \longrightarrow & \QQ^{\times}\\
        & \sigma & \longmapsto & \displaystyle\prod_{1\le i < j \le n} \frac{\sigma(j) - \sigma(i)}{j - i}.
    \end{array}$$
    Let $\Theta$ be a set of all $A \in \mathscr{P}(\{1,\cdots,n\})$ of cardinality $2$.
    For any $\sigma, \tau \in \mathfrak{S}_n$, we have
    $$\begin{aligned}
        \mathrm{sgn}(\tau \circ \sigma) & = \prod_{\{i,j\} \in \Theta} \frac{\tau(\sigma(j)) - \tau(\sigma(i))}{j - i}\\
        & = \prod_{\{i,j\} \in \Theta} \frac{\tau(\sigma(j)) - \tau(\sigma(i))}{\sigma(j) - \sigma(i)} \cdot \prod_{\{i,j\} \in \Theta} \frac{\sigma(j) - \sigma(i)}{j - i}\\
        & = \mathrm{sgn}(\tau) \cdot \mathrm{sgn}(\sigma).
    \end{aligned}$$
    Therefore $\mathrm{sgn}$ is a homomorphism of groups.
    We check the transposition condition.
    Let $\tau$ be of the form $(a\ b)$, $a\neq b$, $a,b \in \{1,\cdots,n\}$
    $$ \mathrm{sgn} (\tau) = \frac{b - a}{a - b} \left(\prod_{i \in \{1, \cdots, n\} \setminus \{a,b\}} \frac{i - a}{i - b}\frac{i - b}{i - a}\right) \left(\prod_{i,j \in \{1,\cdots,n\}\setminus\{a,b\}}\frac{i - j}{i - j}\right) = -1.$$
\end{proofenv}

\begin{propositionenv}\label{4.4.3}
    Let $V$ be a $K$-module and $n$ be a positive integer.
    Let $(x_1,\cdots,x_n) \in V^n$.
    One has
    $$ x_1 \wedge \cdots \wedge x_n = \mathrm{sgn}(\sigma) x_{\sigma(1)} \wedge \cdots \wedge x_{\sigma(n)}$$
    for any $\sigma \in \mathfrak{S}_n$.
\end{propositionenv}
\begin{proofenv}
    The case $n = 1$ is trivial.
    We first treat the case $n = 2$.
    Let $x,y \in V$,
    $$ 0 = \left(x + y\right)\wedge\left(x + y\right) = x\wedge x + x\wedge y + y \wedge x + y \wedge y$$
    so $x \wedge y = - y \wedge x$.

    \quad For the general case, since any $\sigma \in \mathfrak{S}_n$ can be written as the composition of adjacent transpositions, we may assume that $\sigma$ is of the form $(i\; i+1)$ for some $i \in \{1,\cdots, n-1\}$.
    We have
    $$ x_{i+1} \wedge x_i = - x_i \wedge x_{i+1}$$
    so
    $$\begin{aligned}
        & x_1 \wedge \cdots \wedge x_{i-1} \wedge x_{i+1} \wedge x_i \wedge x_{i+2} \wedge \cdots \wedge x_n\\
        = & - x_1 \wedge \cdots \wedge x_{i-1} \wedge x_i \wedge x_{i+1} \wedge x_{i+2} \wedge \cdots \wedge x_n.
    \end{aligned}$$
\end{proofenv}

\begin{remark}
    By proposition \ref{4.4.3}, we obtain that the kernel of the homomorphism:
    $$\begin{array}{rcl}
        V^{\otimes n} & \longrightarrow & \bigwedge^n (V)\\
        x_1 \otimes \cdots \otimes x_n & \longmapsto & x_1 \wedge \cdots \wedge x_n
    \end{array}$$
    is generated by elements of the form $x_1 \otimes \cdots \otimes x_n$ where at least two elements $x_i$, $x_j$ are equal.
\end{remark}

\begin{definitionenv}
    Let $V$, $M$ be $K$-modules and $n$ be a positive integer.
    Let $f: V^n \rightarrow M$ be a $n$-linear mapping.
    If for $(x_1,\cdots,x_n) \in V^n$ such that $x_i = x_{i+1}$ for some $i \in \{1,\cdots,n-1\}$, one has $f(x_1,\cdots,x_n) = 0$, then $f$ is called an \textbf{alternating} $n$-linear mapping.
\end{definitionenv}

\begin{propositionenv}
    Let $V$, $M$ be $K$-modules, $n\in \NN_{\ge 1}$ and $f: V^n \rightarrow M$ be an alternating $n$-linear mapping.
    There exists a unique $K$-linear mapping $\hat{f} : \bigwedge^n (V) \rightarrow M$ such that for any $\left(x_1,\cdots, x_n\right) \in V^n$, 
    $$\hat{f}\left(x_1 \wedge \cdots \wedge x_n\right) = f\left(x_1,\cdots,x_n\right).$$
    \begin{center}
        \begin{tikzcd}
        & V^n \arrow[d, ""] \arrow[dr, "f"] \arrow[dl,""]& \\
        V\otimes \cdots \otimes V \arrow[r, ""] & \bigwedge^n (V) \arrow[dashed, r, swap, "\exists! \hat{f}"]& M
        \end{tikzcd}
    \end{center}
\end{propositionenv}
\begin{proofenv}
    Uniqueness is by definition.
    We prove the existence.
    By theorem \ref{thm:4.1.1}, there exists a unique $K$-linear mapping $\tilde{f} : V^{\otimes n} \rightarrow M$ such that for any $(x_1,\cdots,x_n) \in V^n$,
    $$ \tilde{f}\left(x_1 \otimes \cdots \otimes x_n\right) = f\left(x_1,\cdots,x_n\right).$$
    %\begin{center}
    %    \begin{tikzcd}
    %        V^n \arrow[d, ""] \arrow[dr, "f"]& \text{\small universal property of }\otimes\\
    %        V^{\otimes n} \arrow[dashed, r, swap, "\exists! \tilde{f}"]& M
    %    \end{tikzcd}
    %\end{center}
    Therefore $\tilde{f}$ induces by quotient a $K$-linear mapping $\hat{f} : \bigwedge^n (V) \rightarrow M$ such that for any $(x_1,\cdots,x_n) \in V^n$,
    $$ \hat{f}\left(x_1 \wedge \cdots \wedge x_n\right) = \tilde{f}\left(x_1 \otimes \cdots \otimes x_n\right) = f\left(x_1,\cdots,x_n\right).$$
\end{proofenv}

\begin{exampleenv}[\textbf{Important!}]
    Let $V$ be a $K$-module and $n\in \NN_{\ge 1}$.
    Let $f_1,\cdots,f_n : V \rightarrow K$ be $K$-linear forms of $V$.
    Consider the $n$-linear mapping
    $$\begin{array}{rcl}
        V^n & \longrightarrow & K\\
        (x_1,\cdots,x_n) & \longmapsto & \displaystyle\sum_{\sigma \in \mathfrak{S}_n} \mathrm{sgn}(\sigma) f_{\sigma(1)}(x_1) \cdots f_{\sigma(n)}(x_n).
    \end{array}$$
    If $i,j \in \{1,\cdots,n\}$ are such that $i<j$ and $x = x_i = x_j$, $x\in V$, then
    $$\begin{aligned}
        & \displaystyle\sum_{\sigma \in \mathfrak{S}_n} \mathrm{sgn}(\sigma) f_{\sigma(1)}(x_1) \cdots f_{\sigma(n)}(x_n)\\
        =& - \displaystyle\sum_{k<l} \left(\sum_{\substack{\sigma \in \mathfrak{S}_n\\ \sigma(k) = i, \sigma(l) = j}} \mathrm{sgn}(\sigma) f_{\sigma(1)}(x_1) \cdots f_{\sigma(n)}(x_n)\right. \\
        & \left.+ \sum_{\substack{\sigma \in \mathfrak{S}_n\\ \sigma(k) = j, \sigma(l) = i}} \mathrm{sgn}(\sigma) f_{\sigma(1)}(x_1) \cdots f_{\sigma(n)}(x_n)\right)\\
        =& 0
    \end{aligned}$$
    Therefore defines an alternating linear form on $V^n$, and by universal property of tensor product induces also an alternating linear form on $V^{\otimes n}$.
    By passing to the quotient, it induced a $K$-linear form on $\bigwedge^n (V)$ which we denote by $f_1 \wedge \cdots \wedge f_n$. i.e.
    $$ f_1 \wedge \cdots \wedge f_n : \bigwedge^n (V) \rightarrow K$$
    is $K$-linear.
    
    \quad Let $V$ be a free $K$-module of finite type with a basis, i.e.
    $$ e_{i}^\vee : V\rightarrow K,\; e_{i}^\vee(a_1 e_1 + \cdots + a_n e_n) = a_i.$$
    The following equality holds:
    $$ \left(e_1^\vee \wedge \cdots \wedge e_n^\vee\right) \left(e_1 \wedge \cdots \wedge e_n\right) = \sum_{\sigma \in \mathfrak{S}_n}\mathrm{sgn}(\sigma)e_{\sigma(1)}^\vee(e_1) \cdot \cdots \cdot e_{\sigma(n)}^\vee(e_n) = 1.$$
\end{exampleenv}


\section{Linear Independence Criterion}

\begin{propositionenv}\label{prop:4.5.1}
    Let $V$ be a $K$-module and $n\in \NN_{\ge 1}$.
    Let $(e_i)_{i=1}^n$ be a family of elements of $V$.
    For any $i \in \{1,\cdots,n\}$, let $\left(a_{i,1}, \cdots, a_{i,n}\right)_{i=1}^{n} \in K^n$
    and $x_i = a_{i,1} e_1 + \cdots + a_{i,n} e_n$.
    Then the following equality holds:
    $$ x_1 \wedge \cdots \wedge x_n = \left(\sum_{\sigma \in \mathfrak{S}_n} \mathrm{sgn}(\sigma) a_{1,\sigma(1)} \cdots a_{n,\sigma(n)}\right)e_1 \wedge \cdots \wedge e_n.$$
\end{propositionenv}
\begin{proofenv}
    By multilinearity of wedge product, 
    $$\begin{aligned}
        & x_1 \wedge \cdots \wedge x_n = \sum_{j_1,\cdots,j_n \in \{1,\cdots,n\}} a_{1,j_1} \cdots a_{n,j_n} e_{j_1} \wedge \cdots \wedge e_{j_n}\\
        =& \sum_{\sigma \in \mathfrak{S}_n} a_{1,\sigma(1)} \cdots a_{n,\sigma(n)} e_{\sigma(1)} \cdots e_{\sigma(n)}
    \end{aligned}$$
    where the second equality comes from the fact that if there are distinct indices $k$ and $l$ such that $j_{k} = j_{l}$, then $e_{j_1} \wedge \cdots \wedge e_{j_k} \wedge \cdots \wedge e_{j_l} \wedge \cdots \wedge e_{j_n} = 0$.
    Since $e_{\sigma(1)} \wedge \cdots \wedge e_{\sigma(n)} = \mathrm{sgn} (\sigma) e_1 \wedge \cdots \wedge e_n$, we obtain the equality.
\end{proofenv}

\begin{theoremenv}\label{thm:4.5.2}
    Let $V$ be a $K$-module and $n\in \NN_{\ge 1}$.
    Let $\left(e_i\right)_{i=1}^{n}$ be a family of elements in $V$.
    \begin{enumerate}
        \item Assume that, for any $\lambda \in K\setminus\{0\}$, one has
        $$\lambda e_1 \wedge \cdots \wedge e_n \neq 0,$$
        then $\left(e_i\right)_{i=1}^{n}$ is $K$-linearly independent.
        \item Assume that $\left(e_i\right)_{i=1}^{n}$ forms a system of generators of $V$.
        For any $d \in \NN$, the family
        $$ \left(\lambda e_{i_1} \wedge \cdots \wedge e_{i_d}\right)_{1\le i_1< \cdots < i_d \le n}$$
        generates $\bigwedge^d (V)$.
        In particular, $\bigwedge^d (V) = 0$ if $d>n$.
        \item If $V$ is a free $K$-module and $\left(e_i\right)_{i=1}^{n}$ is a basis of $V$,
        then for any $d\in \NN$, $\left(e_{i_1} \wedge \cdots \wedge e_{i_d}\right)_{1\le i_1 < \cdots < i_d \le n}$ forms a basis of $\bigwedge^d (V)$.
    \end{enumerate}
\end{theoremenv}
\begin{proofenv}
    \quad
    \begin{enumerate}
        \item Assume that $\left(e_i\right)_{i=1}^{n}$ is $K$-linearly dependent.
        Without loss of generality, we may assume that there exists $\lambda \in K \setminus \{0\}$ such that $\lambda e_1$ is a $K$-linear combination of $e_2, \cdots, e_n$.
        Then one has
        $$ \lambda e_1 \wedge \cdots \wedge e_n = \left(a_2 e_2 + \cdots a_n e_n\right) \wedge e_2 \wedge \cdots \wedge e_n = 0.$$
        \item Let $(x_1,\cdots, x_d) \in V^d$.
        For any $i \in \{1,\cdots,d\}$, one can express $x_i$ as a linear combination of $\left(e_i\right)_{i=1}^n$, i.e.
        $$ x_i = \sum_{j=1}^n a_{ij} e_j.$$
        Therefore, $x_1 \wedge \cdots \wedge x_d$ can be expressed as a linear combination of elements of the form $e_{j_1} \wedge \cdots \wedge e_{j_d}$, where $j_1,\cdots,j_d \in \{1,\cdots,n\}$ are pairwise distinct.
        Hence, we deduce that $x_1 \wedge \cdots \wedge x_d$ can be written as a linear combination of the elements of the form $e_{i_1} \wedge \cdots \wedge e_{i_d}$, $1 \le i_1 < \cdots i_d \le n$.
        \item By 2. $\bigwedge^d (V)$ is generated by $\left(e_{i_1} \wedge \cdots \wedge e_{i_d}\right)_{1\le i_1 < \cdots < i_d \le n}$.
        
        Let $\left(\lambda_{i_1,\cdots , i_d}\right)_{1\le i_1 < \cdots <i_d\le n}$ be a family of elements of $K$ such that
        $$ \alpha \coloneq \sum_{1 \le i_1 < \cdots <i_d \le n} \lambda_{i_1,\cdots, i_d} e_{i_1} \wedge \cdots \wedge e_{i_d} = 0.$$
        Let $\left(e_{i}^\vee\right)_{i=1}^n$ be the dual basis of $\left(e_{i}\right)_{i=1}^n$.
        For any $i_1,\cdots,i_d \in \{1,\cdots,n\}$, where $1\le i_1 < \cdots < i_d\le n$, one has
        $$ \left(e_{i_1}^\vee \wedge \cdots \wedge e_{i_d}^\vee\right) \left(\alpha\right) = \lambda_{i_1,\cdots, i_d} = 0.$$
        Hence $\left(e_{i_1} \wedge \cdots \wedge e_{i_d}\right)_{1\le i_1 < \cdots < i_d \le n}$ is linearly independent.
    \end{enumerate}
\end{proofenv}

\begin{corollaryenv}\label{thm:4.5.3}
    Let $V$ be a free $K$-module of finite rank $n$.
    Let $d \in \{1,\cdots,n\}$ and 
    $$ F: \bigwedge^{n-d} (V) \longrightarrow \bigwedge^n (V)$$
    be a $K$-linear mapping.
    There exists a \textit{unique} element $\alpha\in \bigwedge^d (V)$ such that $F(\beta) = \alpha \wedge \beta$ for any $\beta \in \bigwedge^{n-d} (V)$.
\end{corollaryenv}
\begin{proofenv}
    Let $\left(e_i\right)_{i=1}^{n}$ be a basis of $V$ and
    $$ \eta = e_1 \wedge \cdots \wedge e_n$$
    has $\binom{n}{m}$ elements.
    For any $m\in \{1,\cdots, n\}$, let $S_m$ be the set of subsets of $m$ elements of $\{1,\cdots,n\}$.
    For any $I = \{i_1,\cdots, i_m\}\in S_m$, $i_1<\cdots<i_m$, let $B_I = e_{i_1} \wedge \cdots \wedge e_{i_m}$.
    By theorem \ref{thm:4.5.2},
    $$ \left(\beta_{\{1,\cdots,n\}\setminus J}\right)_{J \in S_{m}} = \left(\beta_I\right)_{I \in S_m}$$
    forms a basis of $\bigwedge^m (V)$.
    For any $\beta \in \wedge^{n-d}$, let $f(\beta)$ be the element of $K$ such that 
    $$ F(\beta) = f(\beta) \cdot\eta.$$
    We let 
    $$ \alpha = \sum_{J \in S_{n-d}} \mathrm{sgn} (J) f(\beta_J) \beta_{\{1,\cdots,n\}\setminus J},$$
    where $\mathrm{sgn}(J) \in \{\pm 1\}$ is such that $\beta_{\{1,\cdots,n\}\setminus J} \wedge \beta_J = \mathrm{sgn}(J) \cdot \eta$.
    For any $I \in S_{n-d}$,
    $$\begin{array}{rl}
        \alpha \wedge \beta_{I} &= \sum_{J \in S_{n-d}} \mathrm{sgn} (J) f(\beta_J) \beta_{\{1,\cdots,n\}\setminus J} \wedge \beta_{I}\\
        & = \mathrm{sgn}(I) f(\beta_I) \beta_{\{1,\cdots,n\}\setminus I} \wedge \beta_I \\
        & = f(\beta_I) n = F(\beta_I).
    \end{array}$$
    By linearity, we obtain that $\alpha\wedge \beta = F(\beta)$ for any $\beta \in \bigwedge^{n-d} (V)$.
    If $\alpha'$ and $\alpha$ are two elements such that 
    $$ \alpha \wedge \beta = \alpha' \wedge \beta = F(\beta),\; \forall \beta \in \bigwedge^{n-d} (V),$$
    then
    $$ \left(\alpha - \alpha'\right) \wedge \beta_{I} = 0,\; \forall I \in S_{n-d}.$$
    Hence $\alpha - \alpha' = 0$.
\end{proofenv}


\section{Determinant}

We fix a unitary commutative ring $K$.

\begin{propositionenv}
    Let $V$ be a free $K$-module, $\dim_{K} V = 1$.
    For any $K$-endomorphism $\varphi: V\rightarrow V$, there exists a unique $\lambda_{\varphi} \in K$ such that
    $$ \forall x\in V,\; \varphi(x) = \lambda_{\varphi} x.$$
\end{propositionenv}
\begin{proofenv}
    Let $x_0$ be a generator,
    $$ \varphi(x_0) = \lambda_{\varphi} x_0,\; \lambda_{\varphi} \in K.$$
    For any $x\in V$, $x = bx_0$, hence 
    $$ \varphi(x) = \varphi(b x_0) = b \lambda_{\varphi} x_0 = \lambda_{\varphi} x.$$
\end{proofenv}

\begin{definitionenv}
    Let $V$ be a free $K$-module of rank $n$.
    Let $\varphi : V \rightarrow V$ be a $K$-linear endomorphism, we denote by $\mathrm{det} (\varphi)$ the unique element of $K$ such that for any $x_1,\cdots,x_n \in V$,
    $$ \varphi(x_1) \wedge \cdots \wedge \varphi(x_n) = \det (\varphi) x_1 \wedge \cdots \wedge x_n.$$
    $V^n\rightarrow \bigwedge^n (V)$ is an $n$-linear alternating mapping.
    By the universal property of $\bigwedge^n (V)$,
    \begin{center}
\begin{tikzcd}
	{V^{n}} & {(x_1,\cdots,x_n)} & \\
	{\bigwedge^{n}(V)} & {\bigwedge^{n}(V)} \\
	{x_1 \wedge \cdots \wedge x_n} && {\varphi(x_1) \wedge \cdots \varphi(x_n)}
	\arrow[from=1-1, to=2-1]
	\arrow[from=1-1, to=2-2]
	\arrow[maps to, from=1-2, to=3-3]
	\arrow[dashed, from=2-1, to=2-2]
	\arrow[maps to, from=3-1, to=3-3]
\end{tikzcd}
    \end{center}
\end{definitionenv}

\begin{propositionenv}\label{thm:LinearMap.det_comp}
    Let $V$ be a free $K$-module of rank $n$, and $\varphi: V\rightarrow V$, $\psi: V\rightarrow V$ be $K$-linear endomorphisms.
    Then
    $$ \det (\psi \circ \varphi) = \det (\psi) \det (\varphi).$$
    Moreover,
    $$ \det (\id_V) =1.$$
\end{propositionenv}
\begin{proofenv}
    For any $x_1,\cdots,x_n \in V,$
    $$\begin{aligned}
        \det (\psi \circ \varphi) \left(x_1 \wedge \cdots \wedge x_n\right) &= \det (\psi) \left(\varphi(x_1) \wedge \cdots \wedge \varphi(x_n)\right)\\
        &= \left(\det (\psi) \det (\varphi)\right) x_1 \wedge \cdots \wedge x_n.
    \end{aligned}$$
    One has 
    $$ \id_V(x_1) \wedge \cdots \wedge \id_V(x_n) = x_1 \wedge \cdots \wedge x_n.$$
    so by definition, $\det (\id_V) = 1$.
\end{proofenv}

\begin{theoremenv}
    Let $V$ be a free $K$-module of rank $n$, and $\varphi : V\rightarrow V$ be a $K$-linear endomorphism.
    Then $\varphi$ is a bijection (i.e. a $K$-linear automorphism) if and only if $\det (\varphi)$ is invertible in $K$.
\end{theoremenv}
\begin{proofenv}
    Suppose that $\varphi$ is invertible, then by \ref{thm:LinearMap.det_comp}, 
    $$ \det (\varphi) \det (\varphi^{-1}) = \det (\id_V) = 1.$$
    So $\det (\varphi^{-1}) = \left(\det (\varphi)\right)^{-1}$ in $K$.

    \quad Conversely, assume that $\det (\varphi)$ is invertible in $K$. Let $(e_i)_{i=1}^n$ be a basis of $V$ over $K$.
    For any $i \in \{1,\cdots,n\}$, $x_i = \varphi(e_i)$.
    By 1. of \ref{thm:4.5.2}, $\left(x_i\right)_{i=1}^{n}$ is $K$-linearly independent.
    It suffices to check that $\left(x_i\right)_{i=1}^{n}$ generate $V$.

    \quad Let $\left(e_i^\vee\right)_{i=1}^{n}$ be the dual basis of $\left(e_i\right)_{i=1}^{n}$.
    Put
    $$ g = \left(\det (\varphi)\right)^{-1} \left(e_1^\vee \cdots e_{n}^{\vee}\right),$$
    then $g(x_1, \cdots, x_n) = 1$.
    We define $f: V^{n+1} \rightarrow V$ as 
    $$ \left(y_0,\cdots, y_n\right) \longmapsto \sum_{i=0}^{n}(-1)^n g(y_0, \cdots, y_{i-1}, y_{i+1}, \cdots, y_n) y_i.$$
    It is an alternating $n$-linear mapping.
    Indeed, if $(y_0,\cdots, y_n)$ is such that $y_i = y_{i+1}$, then 
    $$\begin{aligned}
        f(y_0,\cdots,y_n) = & (-1)^{i} f(y_0,\cdots,y_{i-1}, y_{i+1}, \cdots, y_n) y_i\\
        & + (-1)^{i+1} f(y_0,\cdots,y_{i}, y_{i+2}, \cdots, y_n) y_{i+1} = 0.
    \end{aligned}$$
    \begin{center}
\begin{tikzcd}
	{V^{n+1}} & \\
	{0 = \bigwedge^{n+1} (V)} & V
	\arrow[from=1-1, to=2-1]
	\arrow["f \ n\text{-linear alternating}", from=1-1, to=2-2]
	\arrow[dashed, from=2-1, to=2-2]
\end{tikzcd}
    \end{center}
    Let $x\in V$, then
    $$ x\wedge x_1 \wedge \cdots \wedge x_n = 0.$$
    Hence $f(x,x_1,\cdots,x_n) = 0$, so 
    $$ x + \sum_{i=1}^{n} (-1)^{i} g\left(x,x_1,\cdots,x_{i-1},x_{i+1},\cdots,x_n\right)x_i = 0.$$ 
    So $x$ is a linear combination of $\left(x_i\right)_{i=1}^{n}$.
    We conclude that $\left(x_i\right)_{i=1}^{n}$ generate $V$.
\end{proofenv}

\begin{propositionenv}\label{thm:Matrix.det_apply'}
    Let $n$ be a positive integer and let
    $$ A = 
    \begin{pmatrix}
        a_{1,1} & a_{1,2} & \cdots & a_{1,n}\\
        a_{2,1} & a_{2,2} & \cdots & a_{2,n}\\
        \vdots & \vdots & \ddots & \vdots\\
        a_{n,1} & a_{n,2} & \cdots & a_{n,n}
    \end{pmatrix}
    $$
    Then 
    $$ \det (A) = \sum_{\sigma \in \mathfrak{S}_n} \sgn (\sigma) a_{1,\sigma(1)} a_{2,\sigma(2)} \cdots a_{n,\sigma(n)}.$$
\end{propositionenv}
\begin{proofenv}
    Let $\left(e_i\right)_{i=1}^{n}$ be the canonical basis of $K^n$.
    For any $ A(e_i) = a_{i,1} e_1 + \cdots + a_{i,n} e_n$, by lemma \ref{prop:4.5.1}, we have
    $$ A(e_1) \wedge \cdots \wedge A(e_n) = \left(\sum_{\sigma \in \mathfrak{S}_n} \sgn (\sigma) a_{1,\sigma(1)} \cdots a_{n,\sigma(n)}\right)e_1 \wedge \cdots \wedge e_n.$$
\end{proofenv}

\begin{definitionenv}
    Let $V$ and $W$ be $K$-modules and $\varphi : V\rightarrow W$ be $K$-linear mapping.
    We denote by 
    $$ \varphi^\vee : W^\vee \rightarrow V^\vee$$
    the $K$-linear mapping that sends $\alpha \in W^\vee$ to $\varphi^\vee(\alpha) \coloneq \alpha \circ \varphi$.
\end{definitionenv}

\begin{propositionenv}\label{thm:LinearMap.det_dualMap}
    Let $V$ be a free $K$-module of rank $n$, and $\varphi : V \rightarrow V$ be a $K$-linear endomorphism.
    Then
    $$ \det (\varphi^\vee) = \det (\varphi).$$
\end{propositionenv}
\begin{proofenv}
    Let $\left(e_i\right)_{i=1}^{n}$ be a basis of $V$ and $\left(e_i^{\vee}\right)_{i=1}^{n}$ be the dual basis.
    We want to evaluate 
    $$ \varphi^\vee (e_1) \wedge \cdots \wedge \varphi^\vee (e_n)$$
    for some $\alpha_1,\cdots,\alpha_n \in V^\vee$ in terms of $\alpha_1 \wedge \cdots \wedge \alpha_n \in \bigwedge^n (V^\vee)$.
    Let us take $\alpha_1 = e_1^\vee$, $\alpha_2 = e_2^\vee$, …, $ \alpha_n = e_n^\vee$.
    $$\begin{aligned}
        &\left(e_1^\vee \circ \varphi\right) \wedge \cdots \left(e_n^\vee \circ \varphi\right)(e_1 \wedge \cdots \wedge e_n)\\
        = & \left(e_1^\vee \wedge \cdots \wedge e_n^\vee\right)\left(\varphi(e_1) \wedge \cdots \wedge \varphi(e_n)\right)\\
        = & \det (\varphi) \cdot 1.
    \end{aligned}$$
    So,
    $$\left(e_1^\vee \circ \varphi\right) \wedge \cdots \left(e_n^\vee \circ \varphi\right) = \det (\varphi) \left(e_1^\vee \wedge \cdots \wedge e_n^\vee\right).$$
    We deduce that 
    $$ \det (\varphi^\vee) = \det (\varphi).$$
\end{proofenv}

\begin{corollaryenv}\label{thm:Matrix.det_transpose}
    Let $n$ be a positive integer and let $A \in M_n(K)$ be a $n$-times-$n$ matrix.
    One has
    $$ \det (A^{\mathsf{T}}) = \det (A).$$
\end{corollaryenv}


\section{Adjugate}

In this section, we fix free $K$-module of rank $n$ $E$.

\begin{theoremenv}\label{thm:4.7.1}
    For any $\varphi \in \End_K (E)$, there exists $\varphi^a \in \End_K (E)$ such that for any $x_1,\cdots,x_n \in E,$ one has
    $$ x_1 \wedge \varphi(x_2) \wedge \cdots \wedge \varphi(x_n) = \varphi^a (x_1) \wedge x_2 \wedge \cdots \wedge x_n.$$
    Moreover, 
    $$ \varphi^a \circ \varphi = \det (\varphi) \id_E.$$
\end{theoremenv}
\begin{proofenv}
    For any $x_1 \in E$ the mapping
    $$ \begin{array}{rcl}
        \bigwedge^{n-1} (E) & \longrightarrow & \bigwedge^n (E)\\
        x_2 \wedge \cdots \wedge x_n & \longmapsto & x_1 \wedge \varphi(x_2) \wedge \cdots \wedge \varphi(x_n)
    \end{array}$$
    is $K$-linear.
    \begin{center}
    \begin{tikzcd}
	{\bigwedge^{n-1} (E)} & \\
	{E^{n-1}} & {\bigwedge^{n}(E)} \\
	{(x_2,\cdots,x_n)} & {x_1 \wedge \varphi(x_2)\wedge \cdots \wedge \varphi(n)}
	\arrow[from=1-1, to=2-2]
	\arrow[from=2-1, to=1-1]
	\arrow[from=2-1, to=2-2]
	\arrow[maps to, from=3-1, to=3-2]
    \end{tikzcd}
    \end{center}
    Therefore, by \ref{thm:4.5.3}, there exists a unique element of $E$ that we will denote $\varphi^a(x_1)$, such that 
    $$ \varphi^a(x_1) \wedge x_2 \wedge \cdots \wedge x_n = x_1 \wedge \varphi(x_2) \wedge \cdots \wedge \varphi(x_n).$$
    The uniqueness of $\varphi^a(x_1)$ shows that the mapping
    $$ \begin{array}{rrcl}
        \varphi^a : & E & \longrightarrow & E\\
        & x_1 & \longmapsto & \varphi^a(x_1)
        \end{array}$$
    is $K$-linear.
    $$ \varphi^a(\varphi(x_1)) \wedge x_2 \wedge \cdots \wedge x_n = \varphi(x_1) \wedge \varphi(x_2) \wedge \cdots \wedge \varphi(x_n) = \det (\varphi) x_1 \wedge x_2 \wedge \cdots \wedge x_n.$$
    Hence, $\varphi^a \circ \varphi = \det (\varphi) \id_E$.
\end{proofenv}

\begin{definitionenv}
    Let $n$ be a positive integer and $A$ be a $n$-times-$n$ matrix over $K$.
    For $i,j \in \{1,\cdots,n\}$, let $\hat{A}_{i,j}$ be a $(n-1)$-times-$(n-1)$ matrix obtained from $A$ by deleting the $i$-th row and the $j$-th column and let
    $$ A^a_{i,j} = (-1)^{i+j}\det (\hat{A}_{i,j}).$$
    The matrix
    $$\begin{pmatrix}
        A^a_{1,1} & \cdots & A^a_{1,n}\\
        \vdots & \ddots & \vdots\\
        A^a_{n,1} & \cdots & A^a_{n,n}
    \end{pmatrix}$$
    is called the \textbf{adjugate} of $A$, denoted by $A^a$.

    \quad Note that, if we consider $A$ as an endomorphism of $K$, then $A^a$ identifies with the matrix of the endomorphism constructed in the theorem \ref{thm:4.7.1}
\end{definitionenv}

\begin{corollaryenv}
    Let $n$ be a positive integer and
    $$A = \begin{pmatrix}
        a_{1,1} & \cdots & a_{1,n}\\
        \vdots & \ddots & \vdots\\
        a_{n,1} & \cdots & a_{n,n}
    \end{pmatrix}$$
    be an element of $M_n(K)$.
    Then the following equality holds:
    $$ AA^a = \det (A) I_n.$$
    In particular, for any $i\in \{1,\cdots,n\}$,
    $$ \det (A) = \sum_{i=1}^{j} (-1)^{i+j} a_{i,j} A^a_{i,j}.$$
\end{corollaryenv}
